\chapter{Emissão de 2ª via de CIP}
\label{chap:pf-2-via}

\section{O que é?}

É a emissão de nova CIP devido a alteração de nome, inserção de registro
de especialista, alteração de inscrição provisória para definitiva,
roubo ou extravio do documento original\emph{.}

\section{Quem pode utilizar esse serviço?}

Profissionais graduadas/os em Psicologia que atuem no estado de São Paulo e já tenham inscrição no CRP~SP ativa.

\section{Como solicitar o serviço?}

\href{https://www.crpsp.org/agendamento}{Agendamento prévio} para apresentação dos documentos originais na subsede de referência (confira
os endereços \href{https://www.crpsp.org/sede/index}{no \emph{site}}).

\section{Que informações ou documentos são necessários?}

\begin{itemize}
	\item
	Documento de identificação (RG, CNH etc.);
	\item
	Cadastro de Pessoa Física (CPF);
	\item
	diploma de formação em Psicologia;
	\item
	deferimento de registro de especialista;
	\item
	certidão de casamento ou de união estável;
	\item
	boletim de ocorrência (se a CIP foi roubada ou extraviada).
\end{itemize}

\section{Quais são as etapas para a realização desse serviço?}

\begin{enumerate}
	\item
	A/o psicóloga/o agenda o atendimento e paga a taxa de emissão da CIP;
	\item
	a/o psicóloga/o apresenta os documentos originais na subsede de referência (confira os endereços
	\href{https://www.crpsp.org/sede/index}{no \emph{site}}) e, se necessário, faz a coleta de dados biométricos;
	\item
	o CRP~SP emite a segunda via da CIP.
\end{enumerate}

\section{Quanto custa?}

\begin{itemize}
	\item
	R\$~70,83 (valor de emissão da CIP em 2024).
\end{itemize}

\section{Quanto tempo leva?}

30 dias úteis a partir da data de solicitação presencial e coleta dos dados biométricos.

\section{Legislação relacionada}

\begin{itemize}
	\item
	\href{https://atosoficiais.com.br/cfp/resolucao-administrativa-financeira-n-3-2007-institui-a-consolidacao-das-resolucoes-do-conselho-federal-de-psicologia?origin=instituicao&q=resolu\%C3\%A7\%C3\%A3o\%203/2007}{Resolução	CFP nº~03/2007} (Consolidação das Resoluções do Conselho Federal de
	Psicologia);
	\item
	\href{https://transparencia.cfp.org.br/wp-content/uploads/sites/7/2020/05/Manual-de-Procedimentos-Administrativos-e-Financeiro-CFP-RES-20.2018.pdf}{\href{https://atosoficiais.com.br/cfp/resolucao-administrativa-financeira-n-20-2018-altera-o-manual-de-procedimentos-administrativo-e-financeiro-do-sistema-conselhos-de-psicologia-anexo-da-resolucao-cfp-n-202018-a-resolucao-cfp-n-03-2007-a-resolucao-cfp-n-16-2019-e-da-outras-providencias}{Resolução CFP nº~20/2018}}, que determina a revisão e ampliação do \href{https://transparencia.cfp.org.br/wp-content/uploads/sites/7/2020/05/Manual-de-Procedimentos-Administrativos-e-Financeiro-CFP-RES-20.2018.pdf}{\emph{Manual de procedimentos administrativos e financeiros}};
	\item
	\href{https://atosoficiais.com.br/cfp/resolucao-administrativa-financeira-n-2-2021-dispoe-sobre-a-confeccao-a-expedicao-e-o-recolhimento-de-carteiras-de-identidade-profissional-das-os-psicologas-os-revoga-o-capitulo-vi-da-resolucao-cfp-no-03-2007-os-1o-2o-4o-5o-do-art-2o-da-resolucao-cfp-no-02-2007-1o-4o-5o-do-art-2o-e-1o-do-art-4o-da-resolucao-cfp-no-01-1984-e-da-outras-providencias}{Resolução
		CFP nº~02/2021}, que trata da confecção, a expedição e o recolhimento de carteiras de identidade profissional das/os psicólogas/os.
\end{itemize}

\section{Serviços correlatos}

\begin{itemize}
	\item
	\hyperref[chap:pf-cip]{Entrega da carteira de identidade profissional (CIP)}.
\end{itemize}

\sumario